% © 2025 Ing. David Jaroš — CC BY-NC-ND 4.0
%
% This work is licensed under a Creative Commons Attribution-NonCommercial-NoDerivatives
% 4.0 International License (CC BY-NC-ND 4.0).
%
% License History: Earlier drafts (up to v0.3) were released under CC BY 4.0.
% From v0.4 onward, all material is released under CC BY-NC-ND 4.0 to protect
% the integrity of the theoretical work during ongoing academic development.
%
% See LICENSE.md for full license text.

\ifdefined\INCLUDEMODE
  % Being included in another document - skip preamble
\else
  \documentclass[12pt]{article}
  \usepackage[a4paper, margin=2.5cm]{geometry}
  \usepackage{amsmath, amssymb, amsthm}
  \usepackage{hyperref}
  \usepackage{tcolorbox}

  \newtheorem{theorem}{Theorem}[section]
  \newtheorem{lemma}[theorem]{Lemma}
  \newtheorem{corollary}[theorem]{Corollary}
  \newtheorem{proposition}[theorem]{Proposition}
  \theoremstyle{definition}
  \newtheorem{definition}[theorem]{Definition}
  \theoremstyle{remark}
  \newtheorem{remark}[theorem]{Remark}

  \title{\textbf{GR Closure: Bridge Theorem\\
         Equivalence of Derivative-Based and Tetrad-Based Metric Formulas}}
  \author{Ing. David Jaroš}
  \date{March 2026}

  \begin{document}
  \maketitle

  \begin{abstract}
  The UBT literature contains two metric formulas: a derivative-based formula
  $g_{\mu\nu} = \Re[(\partial_\mu\Theta\cdot\partial_\nu\Theta^\dagger)/\mathcal{N}]$
  (appendix\_FORMAL\_emergent\_metric.tex) and a tetrad-based formula
  $g_{\mu\nu} = \Re(\mathrm{Sc}(E_\mu E_\nu^\dagger))$
  (appendix\_R\_GR\_equivalence.tex).  We prove that these two expressions describe the
  \emph{same geometric object} under the natural identification $E_\mu := \partial_\mu\Theta$
  (with normalization absorbed into $\mathcal{N}$).  The normalization factor $\mathcal{N}$
  plays the same role in both formulas and does not introduce any discrepancy.
  \end{abstract}
\fi

\section{Bridge Theorem: Two Metric Formalisms Are Identical}
\label{sec:metric_bridge}

\begin{tcolorbox}[colback=blue!5!white,colframe=blue!75!black,title=The Gap Being Closed]
\textbf{Two coexisting metric formulas in the UBT repository:}

\begin{itemize}
  \item \textbf{Formula~A} (derivative-based, appendix\_FORMAL\_emergent\_metric.tex):
    \[
      g_{\mu\nu}^{(A)} \;:=\; \Re\!\left[\frac{\partial_\mu\Theta\cdot\partial_\nu\Theta^\dagger}{\mathcal{N}}\right]
    \]
  \item \textbf{Formula~B} (tetrad-based, appendix\_R\_GR\_equivalence.tex):
    \[
      g_{\mu\nu}^{(B)} \;:=\; \Re\!\bigl(\mathrm{Sc}(E_\mu E_\nu^\dagger)\bigr)
    \]
\end{itemize}

\textbf{This section proves:} $g_{\mu\nu}^{(A)} = g_{\mu\nu}^{(B)}$ under $E_\mu := \partial_\mu\Theta$.
The two formulas describe the same object.  No redundancy, no contradiction.
\end{tcolorbox}

% -----------------------------------------------------------------------
\subsection{The Natural Identification}

\begin{definition}[Biquaternionic tetrad from $\Theta$]
\label{def:tetrad_natural}
Given the fundamental biquaternionic field $\Theta(q,\tau)$, define the
\emph{biquaternionic coframe tetrad} by
\begin{equation}
  \boxed{E_\mu \;:=\; \partial_\mu\Theta.}
  \label{eq:tetrad_identification}
\end{equation}
This is the canonical, unique identification: the coframe tetrad vectors are
the coordinate derivatives of the fundamental field.
\end{definition}

\begin{remark}
The identification $E_\mu = \partial_\mu\Theta$ is the natural analogue of the
frame-field (vielbein) identification in standard differential geometry, where
one identifies the tangent frame with the pushforward of coordinate basis vectors.
In UBT, the field $\Theta$ plays the role of a biquaternionic ``embedding map'',
and $\partial_\mu\Theta$ are the corresponding pullback frame vectors.
\end{remark}

% -----------------------------------------------------------------------
\subsection{Main Bridge Theorem}

\begin{theorem}[Metric bridge]
\label{thm:metric_bridge}
Under the identification $E_\mu := \partial_\mu\Theta$ and with $\mathcal{N} = 1$
(unnormalized case), Formula~A and Formula~B coincide:
\begin{equation}
  \Re\!\left[\partial_\mu\Theta\cdot\partial_\nu\Theta^\dagger\right]
  \;=\;
  \Re\!\bigl(\mathrm{Sc}(E_\mu E_\nu^\dagger)\bigr).
  \label{eq:bridge_N1}
\end{equation}
With the normalization factor $\mathcal{N} > 0$:
\begin{equation}
  \Re\!\left[\frac{\partial_\mu\Theta\cdot\partial_\nu\Theta^\dagger}{\mathcal{N}}\right]
  \;=\;
  \Re\!\bigl(\mathrm{Sc}(E_\mu' E_\nu'^\dagger)\bigr),
  \qquad
  E_\mu' \;:=\; \frac{\partial_\mu\Theta}{\sqrt{\mathcal{N}}}.
  \label{eq:bridge_normalized}
\end{equation}
\end{theorem}

\begin{proof}
\textbf{Step 1: Scalar part and matrix product.}
For any $A, B \in \mathbb{B} := \mathbb{C}\otimes\mathbb{H}$, the scalar part
$\mathrm{Sc}(AB^\dagger)$ is defined as $\tfrac{1}{4}\mathrm{Tr}(AB^\dagger)$
in the $2\times 2$ matrix representation, where $\mathrm{Tr}$ is the standard
matrix trace.  This equals the coefficient of the identity basis element
$\mathbf{1}$ in the product $AB^\dagger$.

In the UBT literature (appendix\_A, eq.~(56)), the notation
$\partial_\mu\Theta\cdot\partial_\nu\Theta^\dagger$ denotes exactly
$\mathrm{Sc}(\partial_\mu\Theta\cdot\partial_\nu\Theta^\dagger)$, i.e.\ the scalar part
of the biquaternionic product.  We use this identification throughout.

\textbf{Step 2: Substituting $E_\mu = \partial_\mu\Theta$.}
Substituting into Formula~B:
\begin{align}
  \mathrm{Sc}(E_\mu E_\nu^\dagger)
  \;\Big|_{E_\mu = \partial_\mu\Theta}
  &= \mathrm{Sc}((\partial_\mu\Theta)(\partial_\nu\Theta)^\dagger)
   = \mathrm{Sc}(\partial_\mu\Theta\cdot\partial_\nu\Theta^\dagger).
\end{align}
Here $(\partial_\nu\Theta)^\dagger = \partial_\nu\Theta^\dagger$ because the partial
derivative commutes with the biquaternionic adjoint for real coordinates $x^\nu$.

\textbf{Step 3: Taking the real part.}
The real-part operator $\Re$ commutes with the scalar-part operator $\mathrm{Sc}$
(both project to real-valued components via orthogonal projections):
\begin{equation}
  \Re\!\bigl(\mathrm{Sc}(Q)\bigr) = \mathrm{Sc}\!\bigl(\Re(Q)\bigr)
  \quad\forall\; Q\in\mathbb{B}.
\end{equation}
Therefore:
\begin{align}
  \Re\!\bigl(\mathrm{Sc}(E_\mu E_\nu^\dagger)\bigr)
  &= \Re\!\bigl(\mathrm{Sc}(\partial_\mu\Theta\cdot\partial_\nu\Theta^\dagger)\bigr)
   = \Re\!\bigl(\partial_\mu\Theta\cdot\partial_\nu\Theta^\dagger\bigr)
   = g_{\mu\nu}^{(A)}\big|_{\mathcal{N}=1}.
\end{align}

\textbf{Step 4: Including normalization.}
For $E_\mu' = \partial_\mu\Theta/\sqrt{\mathcal{N}}$ with $\mathcal{N}>0$ a real scalar:
\begin{align}
  \mathrm{Sc}(E_\mu' E_\nu'^\dagger)
  &= \mathrm{Sc}\!\!\left(\frac{\partial_\mu\Theta}{\sqrt{\mathcal{N}}}
    \cdot\frac{\partial_\nu\Theta^\dagger}{\sqrt{\mathcal{N}}}\right)
   = \frac{1}{\mathcal{N}}\,\mathrm{Sc}(\partial_\mu\Theta\cdot\partial_\nu\Theta^\dagger).
\end{align}
Taking $\Re$:
\[
  \Re\!\bigl(\mathrm{Sc}(E_\mu' E_\nu'^\dagger)\bigr)
  = \frac{1}{\mathcal{N}}\,\Re\!\bigl(\mathrm{Sc}(\partial_\mu\Theta\cdot\partial_\nu\Theta^\dagger)\bigr)
  = \Re\!\left[\frac{\partial_\mu\Theta\cdot\partial_\nu\Theta^\dagger}{\mathcal{N}}\right]
  = g_{\mu\nu}^{(A)}.
\]
This completes the proof of \eqref{eq:bridge_normalized}.
\end{proof}

% -----------------------------------------------------------------------
\subsection{Interpretation and Consequences}

\begin{tcolorbox}[colback=green!5!white,colframe=green!75!black,title=Verdict: PROVED]
\textbf{The two metric formalisms are identical under $E_\mu = \partial_\mu\Theta$.}

\begin{itemize}
  \item The derivative-based formula (Formula~A) and the tetrad-based formula (Formula~B)
        describe the \emph{same} real symmetric $(0,2)$ tensor.
  \item The identification $E_\mu = \partial_\mu\Theta$ is the unique natural choice.
  \item The normalization factor $\mathcal{N}$ in Formula~A corresponds to normalizing
        the tetrad vectors as $E_\mu' = \partial_\mu\Theta/\sqrt{\mathcal{N}}$ in Formula~B.
  \item Formula~A with $\mathcal{N}=1$ is Formula~B with $E_\mu = \partial_\mu\Theta$.
\end{itemize}
\end{tcolorbox}

\begin{corollary}[Single canonical metric]
\label{cor:canonical}
There is a single canonical UBT emergent metric:
\begin{equation}
  \boxed{g_{\mu\nu}(x)
  \;=\; \Re\!\left[\frac{\partial_\mu\Theta\cdot\partial_\nu\Theta^\dagger}{\mathcal{N}}\right]
  \;=\; \Re\!\bigl(\mathrm{Sc}(E_\mu E_\nu^\dagger)\bigr),
  \quad E_\mu := \frac{\partial_\mu\Theta}{\sqrt{\mathcal{N}}}.}
  \label{eq:canonical_single}
\end{equation}
Both expressions in appendix\_FORMAL\_emergent\_metric.tex and
appendix\_R\_GR\_equivalence.tex refer to this single object.
\end{corollary}

\begin{remark}[Role of $\mathcal{N}$]
The normalization $\mathcal{N}[\Theta]$ is a positive real scalar field that appears
identically in both formalisms — it is part of the \emph{definition} of the tetrad vectors
$E_\mu' = \partial_\mu\Theta/\sqrt{\mathcal{N}}$.  It does \emph{not} introduce any
independent degree of freedom: once $\Theta$ is given, $\mathcal{N}[\Theta]$ is
determined.  Its role is discussed in detail in step3\_signature\_theorem.tex
(Proposition on normalization): $\mathcal{N}$ selects the unit-determinant normalization
of the metric, not its signature.
\end{remark}

\begin{remark}[Equivalence noted in both appendices]
\label{rem:bridge_note}
This bridge theorem should be noted in both appendix files for cross-referencing:
\begin{itemize}
  \item In appendix\_FORMAL\_emergent\_metric.tex: ``The derivative-based and tetrad-based
        metric formulas are equivalent under the identification $E_\mu = \partial_\mu\Theta$;
        see GR\_closure/step1\_metric\_bridge.tex.''
  \item In appendix\_R\_GR\_equivalence.tex: ``The tetrad vectors $E_\mu$ used here
        are identified with $\partial_\mu\Theta$ of the derivative-based formulation;
        see GR\_closure/step1\_metric\_bridge.tex.''
\end{itemize}
\end{remark}

% -----------------------------------------------------------------------
\subsection{Comparison of the Two Approaches}

\begin{center}
\begin{tabular}{lll}
\hline
& \textbf{Formula~A} & \textbf{Formula~B} \\
\hline
Source & appendix\_FORMAL\_emergent\_metric & appendix\_R\_GR\_equivalence \\
Expression & $\Re[\partial_\mu\Theta\cdot\partial_\nu\Theta^\dagger/\mathcal{N}]$ &
            $\Re(\mathrm{Sc}(E_\mu E_\nu^\dagger))$ \\
Tetrad & implicit ($E_\mu' = \partial_\mu\Theta/\sqrt{\mathcal{N}}$) &
         explicit ($E_\mu$ given) \\
Normalization & $\mathcal{N}$ in denominator & absorbed into $E_\mu$ \\
Under $E_\mu=\partial_\mu\Theta$: & \multicolumn{2}{c}{$\bm{g_{\mu\nu}^{(A)} = g_{\mu\nu}^{(B)}}$ \quad (Theorem~\ref{thm:metric_bridge})} \\
\hline
\end{tabular}
\end{center}

% -----------------------------------------------------------------------
\subsection{Summary}

\begin{enumerate}
  \item \textbf{Bridge identified.}
        The natural identification $E_\mu := \partial_\mu\Theta$ (or
        $E_\mu' := \partial_\mu\Theta/\sqrt{\mathcal{N}}$ with normalization)
        directly maps Formula~B to Formula~A.
  \item \textbf{Identity proved.}
        Theorem~\ref{thm:metric_bridge} establishes $g_{\mu\nu}^{(A)} = g_{\mu\nu}^{(B)}$
        as an algebraic identity, not a coincidence or approximation.
  \item \textbf{Single geometric object.}
        The derivative-based and tetrad-based formulas both describe the same emergent
        real symmetric $(0,2)$ tensor field on spacetime.
  \item \textbf{No free parameter.}
        The normalization $\mathcal{N}$ is uniquely determined by $\Theta$ and appears
        consistently in both formulas.
  \item \textbf{Foundation for GR closure.}
        All subsequent GR recovery results (Steps 2--4) may invoke either representation
        of $g_{\mu\nu}$ without ambiguity.
\end{enumerate}

\ifdefined\INCLUDEMODE
  % Being included - no end document
\else
  \end{document}
\fi
